% problem_id: S001-proposition-injective-pic
% source_id: S001 (Stacks Project)
% source_file: algebraization.tex
% statement_locator: algebraization.tex lines 8512--8532
% proof_locator: algebraization.tex lines 8534--8687
% env: proposition
% status: text-extracted-verbatim (difficulty judgement pending, written by hand)
%
% The statement and proof below are the author's own text, extracted verbatim.
% Nothing is paraphrased, shortened, or supplemented.

\begin{proposition}[Koll\'ar]
\label{proposition-injective-pic}
\begin{reference}
\cite[Theorem 1.9]{Kollar-map-pic}
\end{reference}
Let $(A, \mathfrak m)$ be a Noetherian local ring. Let $f \in \mathfrak m$.
Assume
\begin{enumerate}
\item $A$ has a dualizing complex,
\item $f$ is a nonzerodivisor,
\item $\text{depth}(A/fA) \geq 2$, or equivalently $\text{depth}(A) \geq 3$,
\item if $f \in \mathfrak p \subset A$ is a prime ideal with
$\dim(A/\mathfrak p) = 2$, then $\text{depth}(A_\mathfrak p) \geq 2$.
\end{enumerate}
Let $U$, resp.\ $U_0$ be the punctured spectrum of $A$, resp.\ $A/fA$. The map
$$
\Pic(U) \to \Pic(U_0)
$$
is injective. Finally, if (1), (2), (3), $A$ is $(S_2)$, and
$\dim(A) \geq 4$, then (4) holds.
\end{proposition}

\begin{proof}
Let $\mathcal{L}$ be an invertible $\mathcal{O}_U$-module.
Observe that $\mathcal{L}$ maps to $0$ in $\Pic(U_0)$
if and only if we can extend $\mathcal{L}$ to an invertible
coherent triple $(\mathcal{L}, \mathcal{L}_0, \lambda)$
as in Section \ref{section-coherent-triples}.
By Proposition \ref{proposition-hilbert-triple}
the function
$$
n \longmapsto \chi((\mathcal{L}, \mathcal{L}_0, \lambda)^{\otimes n})
$$
is a polynomial $P$. By Lemma \ref{lemma-nonnegative-chi-triple}
we have
$P(n) \geq 0$ for all $n \in \mathbf{Z}$ with equality if and only if
$\mathcal{L}^{\otimes n}$ is trivial. In particular $P(0) = 0$
and $P$ is either identically zero and we win or $P$ has even degree $\geq 2$.

\medskip\noindent
Set $M = \Gamma(U, \mathcal{L})$ and
$M_0 = \Gamma(X_0, \mathcal{L}_0) = \Gamma(U_0, \mathcal{L}_0)$.
Then $M$ is a finite $A$-module of depth $\geq 2$
and $M_0 \cong A/fA$, see proof of Lemma \ref{lemma-nonnegative-chi-triple}.
Note that $H^2_\mathfrak m(M)$ is finite $A$-module by
Local Cohomology, Lemma
\ref{local-cohomology-lemma-local-finiteness-for-finite-locally-free}
and the fact that $H^i_\mathfrak m(A) = 0$ for $i = 0, 1, 2$
since $\text{depth}(A) \geq 3$.
Consider the short exact sequence
$$
0 \to M/fM \to M_0 \to Q \to 0
$$
Lemma \ref{lemma-nonnegative-chi-triple} tells us $Q$ has finite length
equal to $\chi(\mathcal{L}, \mathcal{L}_0, \lambda)$.
We obtain $Q = H^1_\mathfrak m(M/fM)$ and
$H^i_\mathfrak m(M/fM) = H^i_\mathfrak m(M_0) \cong H^i_\mathfrak m(A/fA)$
for $i > 1$ from the long exact sequence of local cohomology
associated to the displayed short exact sequence. Consider the long
exact sequence of local cohomology associated to the sequence
$0 \to M \to M \to M/fM \to 0$. It starts with
$$
0 \to Q \to H^2_\mathfrak m(M) \to H^2_\mathfrak m(M) \to
H^2_\mathfrak m(A/fA)
$$
Using additivity of lengths we see that
$\chi(\mathcal{L}, \mathcal{L}_0, \lambda)$
is equal to the length of the image of
$H^2_\mathfrak m(M) \to H^2_\mathfrak m(A/fA)$.

\medskip\noindent
Let prove the lemma in a special case to elucidate the rest of the proof.
Namely, assume for a moment that $H^2_\mathfrak m(A/fA)$ is
a finite length module. Then
we would have $P(1) \leq \text{length}_A H^2_\mathfrak m(A/fA)$.
The exact same argument applied to $\mathcal{L}^{\otimes n}$ shows that
$P(n) \leq \text{length}_A H^2_\mathfrak m(A/fA)$ for all $n$.
Thus $P$ cannot have positive degree and we win.
In the rest of the proof we will modify this argument to give
a linear upper bound for $P(n)$ which suffices.

\medskip\noindent
Let us study the map
$H^2_\mathfrak m(M) \to H^2_\mathfrak m(M_0) \cong H^2_\mathfrak m(A/fA)$.
Choose a normalized dualizing complex $\omega_A^\bullet$ for $A$.
By local duality
(Dualizing Complexes, Lemma \ref{dualizing-lemma-special-case-local-duality})
this map is Matlis dual to the map
$$
\text{Ext}^{-2}_A(M, \omega_A^\bullet)
\longleftarrow
\text{Ext}^{-2}_A(M_0, \omega_A^\bullet)
$$
whose image therefore has the same (finite) length.
The support (if nonempty) of the finite $A$-module
$\text{Ext}^{-2}_A(M_0, \omega_A^\bullet)$ consists of
$\mathfrak m$ and a finite number of primes
$\mathfrak p_1, \ldots, \mathfrak p_r$ containing $f$ with
$\dim(A/\mathfrak p_i) = 1$. Namely, by
Local Cohomology, Lemma \ref{local-cohomology-lemma-sitting-in-degrees}
the support is contained in the set of primes $\mathfrak p \subset A$ with
$\text{depth}_{A_\mathfrak p}(M_{0, \mathfrak p}) + \dim(A/\mathfrak p) \leq 2$.
Thus it suffices to show there is no prime $\mathfrak p$ containing $f$ with
$\dim(A/\mathfrak p) = 2$ and
$\text{depth}_{A_\mathfrak p}(M_{0, \mathfrak p}) = 0$.
However, because $M_{0, \mathfrak p} \cong (A/fA)_\mathfrak p$
this would give $\text{depth}(A_\mathfrak p) = 1$ which contradicts
assumption (4).
Choose a section $t \in \Gamma(U, \mathcal{L}^{\otimes -1})$
which does not vanish in the points $\mathfrak p_1, \ldots, \mathfrak p_r$, see
Properties, Lemma
\ref{properties-lemma-quasi-affine-invertible-nonvanishing-section}.
Multiplication by $t$ on global sections determines a map $t : M \to A$
which defines an isomorphism $M_{\mathfrak p_i} \to A_{\mathfrak p_i}$ for
$i = 1, \ldots, r$. Denote $t_0 = t|_{U_0}$ the corresponding section
of $\Gamma(U_0, \mathcal{L}_0^{\otimes -1})$ which similarly determines
a map $t_0 : M_0 \to A/fA$ compatible with $t$.
We conclude that there is a commutative diagram
$$
\xymatrix{
\text{Ext}^{-2}_A(M, \omega_A^\bullet) &
\text{Ext}^{-2}_A(M_0, \omega_A^\bullet) \ar[l] \\
\text{Ext}^{-2}_A(A, \omega_A^\bullet) \ar[u]^t &
\text{Ext}^{-2}_A(A/fA, \omega_A^\bullet) \ar[l] \ar[u]_{t_0}
}
$$
It follows that the length of the image of the top horizontal
map is at most the length of $\text{Ext}^{-2}_A(A/fA, \omega_A^\bullet)$
plus the length of the cokernel of $t_0$.

\medskip\noindent
However, if we replace $\mathcal{L}$ by $\mathcal{L}^n$ for $n > 1$,
then we can use
$$
t^n :
M_n = \Gamma(U, \mathcal{L}^{\otimes n})
\longrightarrow
\Gamma(U, \mathcal{O}_U) = A
$$
instead of $t$. This replaces $t_0$ by its $n$th power.
Thus the length of the image of the map
$\text{Ext}^{-2}_A(M_n, \omega_A^\bullet) \leftarrow
\text{Ext}^{-2}_A(M_{n, 0}, \omega_A^\bullet)$
is at most the length of $\text{Ext}^{-2}_A(A/fA, \omega_A^\bullet)$
plus the length of the cokernel of
$$
t_0^n : 
\text{Ext}^{-2}_A(A/fA, \omega_A^\bullet)
\longrightarrow
\text{Ext}^{-2}_A(M_{n, 0}, \omega_A^\bullet)
$$
Via the isomorphism $M_0 \cong A/fA$ the map $t_0$ becomes
$g : A/fA \to A/fA$ for some $g \in A/fA$ and via the corresponding
isomorphisms $M_{n, 0} \cong A/fA$ the map $t_0^n$ becomes
$g^n : A/fA \to A/fA$. Thus the length of the cokernel above
is the length of the quotient of
$\text{Ext}^{-2}_A(A/fA, \omega_A^\bullet)$ by $g^n$.
Since $\text{Ext}^{-2}_A(A/fA, \omega_A^\bullet)$ is a finite $A$-module
with support $T$ of dimension $1$ and since $V(g) \cap T$
consists of the closed point by our choice of $t$
this length grows linearly in $n$ by
Algebra, Lemma \ref{algebra-lemma-support-dimension-d}.

\medskip\noindent
To finish the proof we prove the final assertion. Assume
$f \in \mathfrak m \subset A$ satisfies
(1), (2), (3), $A$ is $(S_2)$, and $\dim(A) \geq 4$.
Condition (1) implies $A$ is catenary, see
Dualizing Complexes, Lemma \ref{dualizing-lemma-universally-catenary}.
Then $\Spec(A)$ is equidimensional by Local Cohomology, Lemma
\ref{local-cohomology-lemma-catenary-S2-equidimensional}.
Thus $\dim(A_\mathfrak p) + \dim(A/\mathfrak p) \geq 4$
for every prime $\mathfrak p$ of $A$. Then
$\text{depth}(A_\mathfrak p) \geq \min(2, \dim(A_\mathfrak p))
\geq \min(2, 4 - \dim(A/\mathfrak p))$ and hence (4) holds.
\end{proof}
